\documentclass[12pt,a4paper]{article}

% Paquetes básicos
\usepackage[utf8]{inputenc}
\usepackage[T1]{fontenc}
\usepackage{amsmath,amssymb}
\usepackage{geometry}
\usepackage{graphicx}
\usepackage{xcolor}
\usepackage{booktabs}
\usepackage{array}
\usepackage{fancyhdr}
\usepackage{lastpage}
\usepackage{hyperref}

% Configuración de página
\geometry{left=2.5cm, right=2.5cm, top=2.5cm, bottom=2.5cm}

% Colores
\definecolor{primary}{RGB}{0,102,204}
\hypersetup{colorlinks=true, linkcolor=primary, urlcolor=primary}

% Encabezado
\pagestyle{fancy}
\fancyhf{}
\fancyhead[L]{\leftmark}
\fancyhead[R]{\thepage/\pageref{LastPage}}
\renewcommand{\headrulewidth}{0.4pt}

% Comandos personalizados
\newcommand{\xor}{\oplus}
\newcommand{\xnor}{\odot}

\begin{document}

% Portada
\begin{titlepage}
\begin{center}

\vspace*{2cm}

{\Huge\bfseries Arquitectura de Computadoras}

\vspace{1cm}

{\Large\itshape Soluciones de Practicas}

\vspace{2cm}

\begin{tabular}{rl}
\textbf{Materia:} & Arquitectura de Computadoras \\
\textbf{Cuatrimestre:} & 2026 \\
\textbf{Fecha:} & \today
\end{tabular}

\vfill

{\large Basado en el libro: Stallings - Arquitectura y Organizacion de Computadores}

\vspace{2cm}

\end{center}
\end{titlepage}

% Indice
\tableofcontents
\newpage

% Contenido
\section{Introduccion}
Este documento contiene las soluciones detalladas de las practicas de Arquitectura de Computadoras.

\section{Practica 1: Representacion de la Informacion}

\subsection{Conversion de bases numericas}

\subsubsection{Ejercicio 1.1: $(25,75)_{10}$ a base 2}

\begin{minipage}[t]{0.48\textwidth}
\textbf{Parte entera (25):}
\begin{verbatim}
25 ÷ 2 = 12 R 1 ↑
12 ÷ 2 = 6  R 0 |
6  ÷ 2 = 3  R 0 |
3  ÷ 2 = 1  R 1 |
1  ↓ 2 = 0  R 1 |
         = 11001
\end{verbatim}
\end{minipage}
\hfill
\begin{minipage}[t]{0.48\textwidth}
\textbf{Parte fraccionaria (0,75):}
\begin{verbatim}
0,75 × 2 = 1,50 → 1
0,50 × 2 = 1,00 → 1
         = 0,11
\end{verbatim}
\end{minipage}

\vspace{0.3cm}
\textbf{Resultado:} $(25,75)_{10} = (11001,11)_2$

\subsubsection{Ejercicio 1.2: $(14,16)_{10}$ a base 2}

\begin{minipage}[t]{0.48\textwidth}
\textbf{Parte entera (14):}
\begin{verbatim}
14 ÷ 2 = 7 R 0 ↑
7  ÷ 2 = 3 R 1 |
3  ÷ 2 = 1 R 1 |
1  ↓ 2 = 0 R 1 |
         = 1110
\end{verbatim}
\end{minipage}
\hfill
\begin{minipage}[t]{0.48\textwidth}
\textbf{Parte fraccionaria (0,16):}
\begin{verbatim}
0,16 × 2 = 0,32 → 0
0,32 × 2 = 0,64 → 0
0,64 × 2 = 1,28 → 1
0,28 × 2 = 0,56 → 0
0,56 × 2 = 1,12 → 1
           0,00101
\end{verbatim}
\end{minipage}

\vspace{0.3cm}
\textbf{Resultado:} $(14,16)_{10} \approx (1110,00101)_2$ (periódico)

\subsubsection{Ejercicio 1.3: $(7,90)_{10}$ a base 2}

\begin{minipage}[t]{0.48\textwidth}
\textbf{Parte entera (7):}
\begin{verbatim}
7 ÷ 2 = 3 R 1 ↑
3 ÷ 2 = 1 R 1 |
1 ↓ 2 = 0 R 1 |
        = 111
\end{verbatim}
\end{minipage}
\hfill
\begin{minipage}[t]{0.48\textwidth}
\textbf{Parte fraccionaria (0,90):}
\begin{verbatim}
0,90 × 2 = 1,80 → 1
0,80 × 2 = 1,60 → 1
0,60 × 2 = 1,20 → 1
0,20 × 2 = 0,40 → 0
0,40 × 2 = 0,80 → 0 (ciclo)
           0,11100
\end{verbatim}
\end{minipage}

\vspace{0.3cm}
\textbf{Resultado:} $(7,90)_{10} \approx (111,11100)_2$ (periódico)

\subsubsection{Ejercicio 1.4: $(48,25)_{10}$ a base 2}

\begin{minipage}[t]{0.48\textwidth}
\textbf{Parte entera (48):}
\begin{verbatim}
48 ÷ 2 = 24 R 0 ↑
24 ÷ 2 = 12 R 0 |
12 ÷ 2 = 6  R 0 |
6  ÷ 2 = 3  R 0 |
3  ÷ 2 = 1  R 1 |
1  ↓ 2 = 0  R 1 |
         = 110000
\end{verbatim}
\end{minipage}
\hfill
\begin{minipage}[t]{0.48\textwidth}
\textbf{Parte fraccionaria (0,25):}
\begin{verbatim}
0,25 × 2 = 0,50 → 0
0,50 × 2 = 1,00 → 1
         = 0,01
\end{verbatim}
\end{minipage}

\vspace{0.3cm}
\textbf{Resultado:} $(48,25)_{10} = (110000,01)_2$

\subsubsection{Ejercicio 1.5: $(23,125)_{10}$ a base 2}

\begin{minipage}[t]{0.48\textwidth}
\textbf{Parte entera (23):}
\begin{verbatim}
23 ÷ 2 = 11 R 1 ↑
11 ÷ 2 = 5  R 1 |
5  ÷ 2 = 2  R 1 |
2  ÷ 2 = 1  R 0 |
1  ↓ 2 = 0  R 1 |
         = 10111
\end{verbatim}
\end{minipage}
\hfill
\begin{minipage}[t]{0.48\textwidth}
\textbf{Parte fraccionaria (0,125):}
\begin{verbatim}
0,125 × 2 = 0,250 → 0
0,250 × 2 = 0,500 → 0
0,500 × 2 = 1,000 → 1
          = 0,001
\end{verbatim}
\end{minipage}

\vspace{0.3cm}
\textbf{Resultado:} $(23,125)_{10} = (10111,001)_2$

\subsubsection{Ejercicio 1.6: $(51,70)_{10}$ a base 2}

\begin{minipage}[t]{0.48\textwidth}
\textbf{Parte entera (51):}
\begin{verbatim}
51 ÷ 2 = 25 R 1 ↑
25 ÷ 2 = 12 R 1 |
12 ÷ 2 = 6  R 0 |
6  ÷ 2 = 3  R 0 |
3  ÷ 2 = 1  R 1 |
1  ↓ 2 = 0  R 1 |
         = 110011
\end{verbatim}
\end{minipage}
\hfill
\begin{minipage}[t]{0.48\textwidth}
\textbf{Parte fraccionaria (0,70):}
\begin{verbatim}
0,70 × 2 = 1,40 → 1
0,40 × 2 = 0,80 → 0
0,80 × 2 = 1,60 → 1
0,60 × 2 = 1,20 → 1
0,20 × 2 = 0,40 → 0
           0,10110
\end{verbatim}
\end{minipage}

\vspace{0.3cm}
\textbf{Resultado:} $(51,70)_{10} \approx (110011,10110)_2$ (periodico)

\subsection{Numeros enteros signados}

Rango C2 con $n$ bits: $[-2^{n-1}, 2^{n-1}-1]$

\subsubsection{Ejercicio 2.1: $(23)_{10}$}

6 bits (rango: $[-32, 31]$). $(23)_{10} = (010111)_2$

\begin{tabular}{ll}
Signo-Magnitud & $0|10111$ \\
C1 & $010111$ \\
C2 & $010111$ \\
Exceso a 32 & $010111 + 100000 = 110111$
\end{tabular}

\subsubsection{Ejercicio 2.2: $(47)_{10}$}

7 bits (rango: $[-64, 63]$). $(47)_{10} = (101111)_2$

\begin{tabular}{ll}
Signo-Magnitud & $0|101111$ \\
C1 & $0101111$ \\
C2 & $0101111$ \\
Exceso a 64 & $0101111 + 1000000 = 1101111$
\end{tabular}

\subsubsection{Ejercicio 2.3: $(-14)_{10}$}

5 bits (rango: $[-16, 15]$). $(14)_{10} = (1110)_2$

\begin{tabular}{ll}
Signo-Magnitud & $1|1110$ \\
C1 & $10001$ (invertir 01110) \\
C2 & $10010$ (C1 + 1) \\
Exceso a 16 & $-14 + 16 = 2 = 00010$
\end{tabular}

\subsubsection{Ejercicio 2.4: $(-21)_{10}$}

6 bits (rango: $[-32, 31]$). $(21)_{10} = (10101)_2$

\begin{tabular}{ll}
Signo-Magnitud & $1|10101$ \\
C1 & $101010$ (invertir 010101) \\
C2 & $101011$ (C1 + 1) \\
Exceso a 32 & $-21 + 32 = 11 = 001011$
\end{tabular}

\subsubsection{Ejercicio 2.5: $(-27)_{10}$}

6 bits (rango: $[-32, 31]$). $(27)_{10} = (11011)_2$

\begin{tabular}{ll}
Signo-Magnitud & $1|11011$ \\
C1 & $100100$ (invertir 011011) \\
C2 & $100101$ (C1 + 1) \\
Exceso a 32 & $-27 + 32 = 5 = 000101$
\end{tabular}

\subsubsection{Ejercicio 2.6: $(213)_{10}$}

9 bits (rango: $[-256, 255]$). $(213)_{10} = (11010101)_2$

\begin{tabular}{ll}
Signo-Magnitud & $0|11010101$ \\
C1 & $011010101$ \\
C2 & $011010101$ \\
Exceso a 256 & $011010101 + 100000000 = 111010101$
\end{tabular}

\subsection{Punto flotante IEEE 754}

Formato: 1 bit signo, 8 bits exponente (exceso 127), 23 bits mantisa

\subsubsection{Ejercicio 3.1: $-1,010_2 \times 2^{-10}$}

Signo: 1 (negativo)\\
Exponente: $-10 + 127 = 117 = 01110101_2$\\
Mantisa: $010$ (sin el 1 implicito)

\textbf{Resultado:} 1 01110101 01000000000000000000000

\subsubsection{Ejercicio 3.2: $0,001_2 \times 2^{-1}$}

Normalizar: $1,0_2 \times 2^{-3} \times 2^{-1} = 1,0_2 \times 2^{-4}$\\
Signo: 0\\
Exponente: $-4 + 127 = 123 = 01111011_2$\\
Mantisa: 0 (todos ceros)

\textbf{Resultado:} 0 01111011 00000000000000000000000

\subsubsection{Ejercicio 3.3: $1_2 \times 2^{-130}$}

Exponente $-130 < -126$ (denormalizado). Desplazar mantisa 4 posiciones:

$1,0_2 \times 2^{-130} = 1,0_2 \times 2^{-4} \times 2^{-126} = 0,0001_2 \times 2^{-126}$

Signo: 0\\
Exponente: 00000000 (denormalizado)\\
Mantisa: $0001$ (sin 1 implicito)

\textbf{Resultado:} 0 00000000 00010000000000000000000

\subsubsection{Ejercicio 3.4: IEEE 754 a notacion cientifica}

Dato: 0 11101011 0100...

Signo: 0 (positivo)\\
Exponente: $11101011_2 = 235_{10} \rightarrow 235 - 127 = 108$\\
Mantisa: $1,0100_2$ (con 1 implicito)

\textbf{Resultado:} $+1,0100_2 \times 2^{108}$

\subsubsection{Ejercicio 3.5: IEEE 754 a notacion cientifica}

Dato: 1 00000000 0101...

Signo: 1 (negativo)\\
Exponente: 00000000 (denormalizado)\\
Mantisa: $0,0101_2$ (sin 1 implicito)\\
Exponente real: $-126$

\textbf{Resultado:} $-0,0101_2 \times 2^{-126}$

\subsubsection{Ejercicio 3.6: IEEE 754 a notacion cientifica}

Dato: 0 11111111 1010...

Exponente: 11111111 = 255 (caso especial)\\
Mantisa: 1010... (distinto de cero)

\textbf{Resultado:} NaN (Not a Number)

\subsubsection{Ejercicio 3.7: $-\infty$ a IEEE 754}

Signo: 1 (negativo)\\
Exponente: 11111111 (todo 1s)\\
Mantisa: 00000000000000000000000 (todo 0s)

\textbf{Resultado:} 1 11111111 00000000000000000000000

\subsubsection{Ejercicio 3.8: NaN a IEEE 754}

Exponente: 11111111 (todo 1s)\\
Mantisa: cualquier valor distinto de cero

\textbf{Resultado:} 0 11111111 10000000000000000000000 (ejemplo)

\subsubsection{Ejercicio 3.9: $122 \times 10^3$ (base 10) a IEEE 754}

$122 \times 10^3 = 122000_{10}$

Convertir a binario: $122000 = 11101110010010000_2 = 1,1101110010010000_2 \times 2^{16}$

Signo: 0\\
Exponente: $16 + 127 = 143 = 10001111_2$\\
Mantisa: $11011100100100000000000$

\textbf{Resultado:} 0 10001111 11011100100100000000000

\subsubsection{Ejercicio 3.10: $-25 \times 10^{-2}$ (base 10) a IEEE 754}

$-25 \times 10^{-2} = -0,25_{10} = -0,01_2 = -1,0_2 \times 2^{-2}$

Signo: 1\\
Exponente: $-2 + 127 = 125 = 01111101_2$\\
Mantisa: 0

\textbf{Resultado:} 1 01111101 00000000000000000000000

\subsubsection{Ejercicio 3.11: $1,01_2 \times 2^{-130}$}

Exponente $-130 < -126$ (denormalizado). Desplazar 4 posiciones:

$1,01_2 \times 2^{-4} \times 2^{-126} = 0,000101_2 \times 2^{-126}$

Signo: 0\\
Exponente: 00000000\\
Mantisa: $000101$

\textbf{Resultado:} 0 00000000 00010100000000000000000

\subsubsection{Ejercicio 3.12: $-11,101_2$ a IEEE 754}

$-11,101_2 = -1,1101_2 \times 2^1$

Signo: 1\\
Exponente: $1 + 127 = 128 = 10000000_2$\\
Mantisa: $1101$

\textbf{Resultado:} 1 10000000 11010000000000000000000

\subsubsection{Ejercicio 3.13: $(23,125)_{10}$ a IEEE 754}

$(23,125)_{10} = 10111,001_2 = 1,0111001_2 \times 2^4$

Signo: 0\\
Exponente: $4 + 127 = 131 = 10000011_2$\\
Mantisa: $0111001$

\textbf{Resultado:} 0 10000011 01110010000000000000000

\subsubsection{Ejercicio 3.14: $-1,101_2 \times 2^{41}$}

Signo: 1\\
Exponente: $41 + 127 = 168 = 10101000_2$\\
Mantisa: $101$

\textbf{Resultado:} 1 10101000 10100000000000000000000

\subsection{Aritmetica en C2 (5 bits)}

Rango: $[-16, 15]$. Overflow: dos positivos dan negativo o dos negativos dan positivo.

\begin{minipage}[t]{0.48\textwidth}
\subsubsection{Ejercicio 4.1: $01111 + 01000$}
\begin{verbatim}
  01111 (+15)
+ 01000 (+8)
---------
  10111 (-9)
\end{verbatim}
\textbf{Resultado:} Overflow
\end{minipage}
\hfill
\begin{minipage}[t]{0.48\textwidth}
\subsubsection{Ejercicio 4.2: $01001 + 11000$}
\begin{verbatim}
  01001 (+9)
+ 11000 (-8)
---------
 100001
\end{verbatim}
Descartar carry: $00001 = +1$\\
\textbf{Resultado:} $+1$
\end{minipage}

\vspace{0.5cm}

\begin{minipage}[t]{0.48\textwidth}
\subsubsection{Ejercicio 4.3: $01000 - 00011$}
Resta = suma C2: $01000 + 11101$
\begin{verbatim}
  01000 (+8)
+ 11101 (-3)
---------
 100101
\end{verbatim}
Resultado: $00101 = +5$\\
\textbf{Resultado:} $+5$
\end{minipage}
\hfill
\begin{minipage}[t]{0.48\textwidth}
\subsubsection{Ejercicio 4.4: $00111 - 10111$}
$00111 + 01001$ (C2 de 10111)
\begin{verbatim}
  00111 (+7)
+ 01001 (+9)
---------
  10000 (-16)
\end{verbatim}
\textbf{Resultado:} Overflow
\end{minipage}

\vspace{0.5cm}

\begin{minipage}[t]{0.48\textwidth}
\subsubsection{Ejercicio 4.5: $11000 + 11101$}
\begin{verbatim}
  11000 (-8)
+ 11101 (-3)
---------
 110101
\end{verbatim}
Resultado: $10101 = -11$\\
\textbf{Resultado:} $-11$
\end{minipage}
\hfill
\begin{minipage}[t]{0.48\textwidth}
\subsubsection{Ejercicio 4.6: $00111 + 00011$}
\begin{verbatim}
  00111 (+7)
+ 00011 (+3)
---------
  01010 (+10)
\end{verbatim}
\textbf{Resultado:} $+10$
\end{minipage}

\vspace{0.5cm}

\begin{minipage}[t]{0.48\textwidth}
\subsubsection{Ejercicio 4.7: $10110 + 10111$}
\begin{verbatim}
  10110 (-10)
+ 10111 (-9)
---------
 101101
\end{verbatim}
Resultado: $01101 = +13$\\
\textbf{Resultado:} Overflow
\end{minipage}
\hfill
\begin{minipage}[t]{0.48\textwidth}
\subsubsection{Ejercicio 4.8: $11110 + 11101$}
\begin{verbatim}
  11110 (-2)
+ 11101 (-3)
---------
 111011
\end{verbatim}
Resultado: $11011 = -5$\\
\textbf{Resultado:} $-5$
\end{minipage}

\vspace{0.5cm}

\begin{minipage}[t]{0.48\textwidth}
\subsubsection{Ejercicio 4.9: $11111 + 01111$}
\begin{verbatim}
  11111 (-1)
+ 01111 (+15)
---------
 101110
\end{verbatim}
Resultado: $01110 = +14$\\
\textbf{Resultado:} $+14$
\end{minipage}
\hfill
\begin{minipage}[t]{0.48\textwidth}
\subsubsection{Ejercicio 4.10: $01101 - 01110$}
$01101 + 10010$ (C2 de 01110)
\begin{verbatim}
  01101 (+13)
+ 10010 (-14)
---------
  11111 (-1)
\end{verbatim}
\textbf{Resultado:} $-1$
\end{minipage}

\vspace{0.5cm}

\begin{minipage}[t]{0.48\textwidth}
\subsubsection{Ejercicio 4.11}
Falta enunciado en el PDF original.
\end{minipage}

\section{Practica 2: Logica Digital}

\subsection{Minimizacion de funciones booleanas}

% Agregar contenido aqui

\subsection{Maquinas de estados finitos}

% Agregar contenido aqui

\end{document}
